
% ============================================================================
\clearpage
\section*{v4.10.0-r15 Source-Only: A Correction of the Frame Reading of v4.9.0-r38, and Status Notes on Shadow, Spin-1 Triangle and Complementarity}
\addcontentsline{toc}{section}{v4.10.0-r15 Source-Only: A Correction of the Frame Reading of v4.9.0-r38, and Status Notes on Shadow, Spin-1 Triangle and Complementarity}
\label{sec:v4100-r15-r38-correction-complementarity}
% ============================================================================

\paragraph{Purpose.}
At the author's request (2026-11-15) layer v4.9.0-r38 was read again against the later frame layers.  Its reading of the
frames as triality slots predates v4.9.0-r51, and one of its statements contradicts r51.1(a); r51 corrected the frame rows
of r41, r43 and r47--r49 but not those of r38.  This layer records that correction.  It also records four explorations of
15 November 2026 that followed v4.10.0-r14: the core read as the shadow of a horizon, the spin-1 triangle, the Koide value
as complete complementarity of a qubit, and the generations read as eigenstates of the three frame maps.  All items are
explorations without pre-registration and without comparison with data; the Koide value \(K=\tfrac23\), the number three
and the phases were known in advance.  The notes are
\texttt{notes/\allowbreak research/\allowbreak v4\_10\_0\_r15/} and the folders cited there.  Earlier layers are not
rewritten.

\noindent\textbf{Source-only status.}
This v4.10.0-r15 layer is source-only.  It does not build a PDF, does not assign a new DOI,
and does not authorize a public release.  The v5.1.0 Version DOI
\href{https://doi.org/10.5281/zenodo.24237112}{10.5281/zenodo.24237112} and the Concept DOI
\href{https://doi.org/10.5281/zenodo.19210728}{10.5281/zenodo.19210728} are unchanged by it.

\subsection*{v4.10.0-r15 correction}
\label{subsec:v4100-r15-correction}

\paragraph{Correction v4.10.0-r15.1 (the frame reading of v4.9.0-r38).}
Layer v4.9.0-r38 works on the real slots of \(\OO^3\), explicitly without the \(\mathbb C\) factor, and reads the frames as
slots: the photon as the fixed vector slot, matter and tachyon as the two spinor slots exchanged by \(\rho\).  Since
v4.9.0-r51 the frame maps act on the doublet through the \(\mathbb C\) factor of the Dixon store: the tachyon map
\(z\mapsto iz\), the time reversal \(z\mapsto\bar z\) and \(J_B\).  The state-disk picture of v4.9.0-r04 and r19 places the
frames on the same doublet, matter and tachyon at the two poles of one axis and the photon at the unpolarised centre.
\begin{enumerate}[leftmargin=2em]
\item[(a)] \emph{Unchanged.}  Theoremlet v4.9.0-r38.1(a)--(e) is algebra on \(\mathfrak{so}(8)^3\) and remains valid: \(\rho\) is outer of
order two, its fixed algebra is \(\mathfrak{so}(7)\), and \(\tau\), \(\rho\) generate \(S_3\) on the slots.
\item[(b)] \emph{Corrected.}  Reading item~6 of v4.9.0-r38 and its register row ``chirality: \(M\leftrightarrow T\) as exchange of the two
spinor slots'' contradict v4.9.0-r51.1(a): the tachyon map fixes every sector and hence the chirality; the ports, and with them
the chiralities, are exchanged by the time reversal.  The frame change \(M\leftrightarrow T\) is not a change of chirality.
\item[(c)] \emph{Superseded as a frame statement.}  The assignment of M, P, T to triality slots (reading items~2 and~3 of
v4.9.0-r38) is an earlier picture.  The slot transposition \(\rho\) stays a symmetry of the slots; it is not shown to be a frame
map.  This affects the name, not the mathematics, of later readings: in v4.9.0-r44 (H29) the Brannen form from a
\(\tau\)-symmetric, \(\rho\)-broken coupling stands, but ``breaking of the frame reflection'' is not covered.  The sentence ``the
frame triole is a transposition (v4.9.0-r38)'' in v4.10.0-r14.3(d) holds for \(\rho\) on the slots, not for the frame maps of
v4.9.0-r51.
\item[(d)] \emph{Open item~5 of v4.9.0-r38.}  The conflict ``on one level the three slots cannot carry both frames and generations''
arises only if both sit on the slots.  With the frames on the \(\mathbb C\) factor and the generations on the slots (v4.9.0-r37) or at
the place of spin (v4.10.0-r14.4), the two roles lie on different tensor factors of \(\mathbb C\otimes\HH\otimes\OO\).  The
separation by zoom level of v4.9.0-r39 and r41 remains a result of the model but is no longer needed for this purpose.
\end{enumerate}
Source: \texttt{notes/\allowbreak research/\allowbreak op3\_frame\_map\_generations\_2026\_11\_15/}.

\subsection*{v4.10.0-r15 status notes}
\label{subsec:v4100-r15-status}

\paragraph{Theoremlet v4.10.0-r15.2 (the core as the shadow of a horizon).}
Following the author's idea that one sphere casts a shadow on the other, read the flux-free core of v4.10.0-r14.6 as the
visible shadow of the guarded collapse kernel (remark \texttt{rem:hc-shadow-vs-core}).  A shadow is the missing kind of map
from a cross-section to a solid angle and is flux-free by construction.  Its size is not fixed: the Synge shadow of a
Schwarzschild hole has \(\Omega=\gamma_M\) only for an observer at \(r=17.99\,M\), a distance given by nothing in the inventory; the
distance from the frame offset (\(1-2M/r=1-\beta_{\rm sym}^2\)) gives \(1.537\) sr; the sector reading of CT-001 gives
\(\Omega=\gamma/2\), half of what the spin-\(\tfrac12\) cap needs.  A spin-1 charge does not remove this factor two (r15.3).  Since a
shadow depends on the observer while masses must not, the shadow that fixes the masses has to be read in the rest frame of
matter; this gives the guard ``end in frame M'' of CT-002 a reason, as a reading.

\paragraph{Theoremlet v4.10.0-r15.3 (the spin-1 triangle).}
Three spin-1 coherent states symmetric about a cycle axis with pairwise overlap \(|\langle n_k|n_l\rangle|^2=\cos^4(a/2)=\tfrac12\)
have side angle \(\cos a=\sqrt2-1\), axis angle \(\cos^2\psi=(2\sqrt2-1)/3\) (\(\psi=38.68^\circ\)), \(K=\tfrac23\), Bargmann phase
\[
\varphi=3\arccos\bigl(1-\tfrac1{\sqrt2}\bigr)-\pi=0.679042,
\]
and \(\det G=0.0503>0\): rank three, root ratios \(1:16.24:65.93\).  One \(\mathfrak{su}(2)\) thus carries the Koide value, three
positive masses and a cycle of order three without any core; the remainder to the lepton phase \(\tfrac23\) is \(0.012376\).  A
search over 143 candidate directions fixed beforehand in the script (lattice, \(A_4\), Hopf and phase angles; tolerance
\(\pm0.05^\circ\)) found no native axis; the nearest are \(\tfrac23\) rad (\(-0.48^\circ\)) and the lattice direction \((2,1,0)\)
(\(+0.56^\circ\)).  The direction is not independent: it follows from the overlap \(\tfrac12\) alone.  In every spin, Koide asks for
three states pairwise half-way, in Fubini--Study distance \(\pi/4\), between equal and orthogonal.

\paragraph{Theoremlet v4.10.0-r15.4 (Koide as complete complementarity).}
The octant consists of one eigenstate each of \(\sigma_x\), \(\sigma_y\), \(\sigma_z\), one state from each of the three mutually
unbiased bases of a qubit.  In dimension \(d\) a complete set has \(d+1\) bases with pairwise overlap \(1/d\); one state from each
gives \(N=d+1\) states with
\[
K=\frac{\operatorname{Tr}G^2}{(\operatorname{Tr}G)^2}=\frac2{d+1}=\frac2N
\]
(checked with Wootters--Fields bases for \(d=2,3,5,7\)), and rank at most \(d<N\).  Only \(d=2\) gives three states, with
\(K=\tfrac23\).  Read this way, ``three generations'' and ``Koide \(\tfrac23\)'' are one statement, \(K=2/N\) at \(N=3\).  Two principles
are assumed: one state per complementary basis, and a complete set.

\paragraph{Theoremlet v4.10.0-r15.5 (one generation per frame map).}
On the doublet of v4.9.0-r51 take one eigenstate each of the time reversal \(T\) (\(\sigma_z\)), of \(J_B\) (\(\sigma_x\)) and of the
tachyon map \(R\colon z\mapsto iz\).  As a real map \(R\) is a rotation with \(R^2=-1\); its eigenvalues are \(\pm i\) and its eigenlines
are those of \(\sigma_y\), so the \(\mathbb C\) factor is needed.
\begin{enumerate}[leftmargin=2em]
\item[(a)] The three eigenbases are pairwise unbiased.  All eight choices give \(K=\tfrac23\) and Bargmann phase \(\pm\pi/4\); the sign
is the orientation of the three Bloch vectors.
\item[(b)] The frame change \(R\) fixes its own eigenstate and sends the other two to their partners; the triple turns by \(\pi\) about
the \(R\) axis, and \(K\) and the sign of the phase are kept.  Like the charges in v4.9.0-r51, the Koide data are the same in M and T.
\item[(c)] Complex conjugation of the states, an antiunitary map, flips only the \(R\) eigenstate and reverses the sign of the phase.
\end{enumerate}
Reading: the two principles of r15.4 become one, a generation per frame map, and the number three of generations is the
number of involutions of the doublet.  Whether the generation index sits on the frame doublet, and whether the time reversal
acts on it linearly or antilinearly (v4.9.0-r02 requires an antilinear involutive exchange of matter and tachyon in three
complex dimensions), is not shown.  The bend to the physical masses remains with the core, \(\gamma_M\) and OP1.

\subsection*{v4.10.0-r15 status in the three frames}
\label{subsec:v4100-r15-frames}

\begin{itemize}[leftmargin=2em]
\item M: r15.2 reads the shadow, and hence the core size, in the matter frame; all mass statements are matter-frame readouts.
\item T: by r15.1(b) and r15.5(b) chirality and the Koide data are the same in M and T.
\item P: the photon is the unpolarised centre of the doublet (v4.9.0-r04, r19); otherwise not evaluated.
\end{itemize}

\subsection*{v4.10.0-r15 register entries}
\label{subsec:v4100-r15-register}

\begin{center}
\small
\begin{tabular}{@{}>{\raggedright\arraybackslash}p{0.34\textwidth}>{\raggedright\arraybackslash}p{0.58\textwidth}@{}}
\toprule
statement & status after v4.10.0-r15\\
\midrule
r38.1(a)--(e) & unchanged (r15.1(a))\\
``\(M\leftrightarrow T\) is a change of chirality'' (r38) & corrected: the tachyon map keeps the chirality; time reversal exchanges it (r15.1(b))\\
frames as triality slots (r38) & earlier picture; frame maps act through the \(\mathbb C\) factor (r15.1(c))\\
frames versus generations on the slots (r38, item~5) & conflict absent in its original form; different tensor factors (r15.1(d))\\
core as shadow & reading; size open, frame M (r15.2)\\
spin-1 triangle & exact; reference direction not native (r15.3)\\
Koide as complementarity & \(K=2/(d+1)\); only \(d=2\) gives three states (r15.4)\\
generations per frame map & reading; unbiased, \(K=\tfrac23\), sign of the phase from time reversal (r15.5)\\
CT-001 to CT-004 & unchanged\\
H100 / H102 & pending / pending\\
\bottomrule
\end{tabular}
\end{center}

\subsection*{v4.10.0-r15 ledger and guard update}
\label{subsec:v4100-r15-ledger-guard}

\[
R_{\rm chk}^{v4100\text{-}r15}=(1,1,0)
\]
for (correction and status notes recorded; formulas checked in the cited scripts; generation index on the frame doublet,
overlap \(\tfrac12\) and core size not derived).

The layer spends no guard:
\[
R_{\rm glob}^{(v4100\text{-}r15,+)}=(0,1,1,1,1),
\qquad
\Gamma_{\rm guard}^{v4100\text{-}r15}=(1,1,1,1,1),
\qquad
\beta_4=1/256\ \text{unspent}.
\]
OP~1, OP~3 and OP~4 remain open.  The layer does not imply exact finite \(C5\) zero closure, exact all-readout terminality,
\(C6/T2\), a global photon-clock/RH theorem, or global all-readout closure.
